% ============================================================================
% CAPÍTULO 3: DESARROLLO DE LA PROPUESTA
% ============================================================================
% Este capítulo detalla el desarrollo práctico de la solución propuesta, la
% implementación del software, arquitectura del sistema, base de datos,
% o el núcleo del trabajo práctico o de ingeniería de tu tesis.
% ============================================================================

\chapter{Desarrollo de la Solución}

% --- Introducción al capítulo ---
% En este capítulo se describe de manera detallada la implementación física y lógica de la solución propuesta...

\section{Arquitectura de la solución}

[Presenta aquí el esquema general de tu sistema (patrón arquitectónico como MVC, microservicios, cliente-servidor, etc.). Puedes incluir diagramas de arquitectura en la carpeta `figuras/` e insertarlos con el comando \texttt{\textbackslash includegraphics}.]

\section{Modelo de datos y almacenamiento}

[Detalla el diseño de tu base de datos o estructura de almacenamiento. Incluye diagramas entidad-relación (DER) o esquemas NoSQL y explica las tablas o colecciones principales.]

\section{Implementación detallada de componentes}

[Describe la programación y construcción de los módulos principales del sistema, incluyendo algoritmos clave, consumo de APIs, servicios en la nube o lógica del negocio relevante.]

\subsection{Componente o módulo backend}
[Descripción técnica del desarrollo del lado del servidor, endpoints principales, etc.]

\subsection{Componente o módulo frontend (Interfaz de usuario)}
[Descripción técnica del desarrollo del lado del cliente, diseño de interfaces, etc.]

\section{Casos de uso principales y flujos de trabajo}

[Explica cómo los usuarios interactúan con el sistema mediante flujos clave. Apóyate de diagramas de secuencia o capturas de pantalla de la interfaz para ilustrar los procesos principales.]
